% ch28-student-journey
\section{يوم في حياة النظام: رحلة طالبة كاملة (وكيف يردّ الآخرون)}

كل فصل سابق شرح جزءًا: الطبقات، و\en{CQRS}، والأمان، والأحداث، والدفع،
والذكاء الاصطناعيّ، والزمن الحقيقيّ. هذا الفصل مختلف: لن نضيف مفهومًا جديدًا،
بل نربط كل ما سبق في قصّة واحدة متّصلة. سنتابع طالبة اسمها «سارة» من
لحظة إنشاء حسابها حتى تسأل مساعد الذكاء الاصطناعيّ سؤالًا، وعند كل خطوة نسأل
ثلاثة أسئلة: أيّ \en{command} أو \en{query} يعمل؟ أيّ \en{handler} ينفّذه؟ وماذا
يفعل الدور الآخر (أدمن، أو استشاريّ، أو جهة مانحة) ردًّا على فِعل سارة؟

\begin{idea}[الفكرة: النظام سلسلة أفعال وردود]
كل ما يحدث في \en{ScholarPath} هو أمرٌ يغيّر الحالة (\en{command})، أو استعلامٌ
يقرأها (\en{query}). لكن الأمر لا يعيش وحده: كثيرٌ منه يُطلق حدثًا دوميًّا
(\en{domain event}) بعد نجاحه، وهذا الحدث يوقظ دورًا آخر في الطرف المقابل. رحلة
سارة إذًا ليست خطًّا مستقيمًا، بل حوارٌ متبادَل بين أربعة أدوار.
\end{idea}

\subsection{خريطة اليوم: الرحلة في لمحة}

قبل التفاصيل، هذه خريطة الرحلة كلها. إحدى عشرة محطّة، ملوّنة حسب المرحلة: التسجيل،
ثم الاكتشاف والذكاء الاصطناعيّ، ثم التقديم، ثم الحجز والدفع، ثم الجلسة، ثم المجتمع.
اقرأيها من أعلى لأسفل — كل صندوق يذكر الأمر الذي يعمل والدور الذي يردّ.

\begin{center}
\begin{LTR}
\begin{tikzpicture}[node distance=0.26cm]
  \node[wbox,fill=spLightGreen,draw=spWhy!50] (a1)
    {\small 1 — سارة تُنشئ حسابها وتؤكّد بريدها \en{(ASP.NET Identity)}};
  \node[wbox,below=0.26cm of a1,fill=spLightGreen,draw=spWhy!50] (a2)
    {\small 2 — تُكمل ملفها والـ \en{onboarding}: المستوى والمجالات والدول المفضّلة};
  \node[wbox,below=0.26cm of a2,fill=spLightPurple,draw=spQA!45] (a3)
    {\small 3 — الذكاء الاصطناعيّ يقترح منحًا ويفحص الأهليّة \en{(LocalAiService)}};
  \node[wbox,below=0.26cm of a3,fill=spLightBlue,draw=spPrimary!45] (a4)
    {\small 4 — تكتشف منحة وتقدّم طلبًا \en{(SubmitApplicationCommand)}};
  \node[wbox,below=0.26cm of a4,fill=spLightBlue,draw=spPrimary!45] (a5)
    {\small 5 — الجهة المانحة تراجع وتردّ (قبول أو رفض بسبب) والإشعار يصل للجرس};
  \node[wbox,below=0.26cm of a5,fill=spLightOrange,draw=spSec!45] (a6)
    {\small 6 — تحجز جلسة استشاريّ وتدفع بـ \en{Stripe}: الحجز يُجمِّد المبلغ};
  \node[wbox,below=0.26cm of a6,fill=spLightOrange,draw=spSec!45] (a7)
    {\small 7 — الاستشاريّ يقبل \en{(AcceptBookingCommand)}: القبض يأخذ المبلغ};
  \node[wbox,below=0.26cm of a7,fill=spLightTeal,draw=spAccent!45] (a8)
    {\small 8 — جلسة فيديو \en{ACS} داخل المتصفّح مع تسجيل تلقائيّ};
  \node[wbox,below=0.26cm of a8,fill=spLightTeal,draw=spAccent!45] (a9)
    {\small 9 — سارة تقيّم الجلسة \en{(SubmitConsultantRatingCommand)}};
  \node[wbox,below=0.26cm of a9,fill=spLightYellow,draw=spYellow!70!black] (a10)
    {\small 10 — تنشر سؤالًا في المجتمع ويصل الردّ لحظيًّا عبر \en{SignalR}};
  \node[wbox,below=0.26cm of a10,fill=spLightYellow,draw=spYellow!70!black] (a11)
    {\small 11 — تسأل مساعد الذكاء الاصطناعيّ فيردّ بإجابة مع مصادرها \en{(RAG)}};

  \foreach \f/\t in {a1/a2,a2/a3,a3/a4,a4/a5,a5/a6,a6/a7,a7/a8,a8/a9,a9/a10,a10/a11}
    \draw[arr] (\f) -- (\t);
\end{tikzpicture}
\end{LTR}
\end{center}

الأخضر مرحلة التسجيل، والبنفسجيّ مرحلة الذكاء الاصطناعيّ، والأزرق مرحلة التقديم،
والبرتقاليّ مرحلة الحجز والدفع، والتركوازيّ مرحلة الجلسة، والأصفر مرحلة المجتمع.
لاحظي أنّ خطوات مثل 5 و7 و10 يقوم بها دورٌ غير سارة — هذه هي الردود.

\subsection{المحطّة الأولى: التسجيل وتأكيد البريد}

تبدأ سارة بنموذج تسجيل بسيط: اسم، وبريد، وكلمة مرور. الواجهة ترسل طلبًا إلى
\code{AuthController}، فيُنشئ \en{ASP.NET Identity} صفَّ مستخدم جديدًا بكلمة مرور
مجزّأة (\en{hashed}) لا مخزّنة كنصّ. النظام «مبوَّب» (\en{gated}): لا شيء يعمل قبل
تأكيد البريد. لذا يُرسَل بريد فيه رابط رمزيّ (\en{token})، وحين تضغط عليه سارة
يتأكّد الحساب ويصبح دورها الافتراضيّ \en{Student}.

\begin{why}
لماذا نُبوِّب النظام خلف تأكيد البريد؟ لأنّ حسابات وهميّة تُلوّث التوصيات، وتُغرق
الجهات المانحة بطلبات مزيّفة، وتستهلك ميزانيّة الذكاء الاصطناعيّ. تأكيد البريد
بوّابةٌ رخيصة تُبقي مجتمع المنصّة حقيقيًّا منذ اللحظة الأولى.
\end{why}

\begin{note}
هنا لا يوجد «دور آخر يردّ» بعد — سارة وحدها. لكنّ كل فِعل حسّاس من الآن فصاعدًا
سيُسجَّل في سجلّ التدقيق (\en{audit log}) عبر سِمة \code{[Auditable]} كما رأينا سابقًا.
\end{note}

\subsection{المحطّة الثانية: استكمال الملف والـ onboarding}

بعد الدخول تمرّ سارة بخطوات \en{onboarding} تجمع بيانات ملفها: مستواها الأكاديميّ
(\en{AcademicLevel})، ومجالات دراستها المفضّلة (\en{PreferredFields})، والدول التي
تريد الدراسة فيها (\en{PreferredCountries}). قد تبدو خطوة شكليّة، لكنّها في الحقيقة
وقود الذكاء الاصطناعيّ: بدونها لا يملك المحرّك ما يطابق به المنح.

\begin{keyp}[نقطة مهمة: الملف هو مُدخَل الخوارزميّة]
كل حقل تملؤه سارة الآن سيُقرَأ حرفيًّا في المحطّة التالية: \code{AcademicLevel} يعطي
٣٠ نقطة عند تطابق المستوى، و\code{PreferredFields} حتى ٤٠ نقطة، و\code{PreferredCountries}
حتى ٢٥ نقطة. الملف الفارغ يعني توصيات بصفر — لذا نحثّ الطالب على إكماله.
\end{keyp}

\subsection{المحطّة الثالثة: الذكاء الاصطناعيّ يقترح ويفحص الأهليّة}

الآن يتدخّل الذكاء الاصطناعيّ لأوّل مرّة. حين تفتح سارة لوحتها، ترسل الواجهة
\code{GenerateRecommendationsCommand}. يستقبله \code{GenerateRecommendationsCommandHandler}
الذي يمرّ أوّلًا ببوّابة التكلفة (\code{AiCostGate}) لئلّا يتجاوز حدّ اليوم، ثمّ
ينادي منفذ \code{IAiService}. في بيئتنا الحاليّة يُحقَن المحوّل \code{LocalAiService}:
محرّكٌ حسابيّ بحت بلا نداء لـ \en{GPT}، يعطي نتائج قابلة للتفسير والاختبار.

هذه هي قلب الخوارزميّة — حلقة تحسب نقاط كل منحة مفتوحة مقابل ملف سارة:

\begin{english}
\begin{lstlisting}[language={[Sharp]C}]
foreach (var c in candidates)
{
    var score = 0;
    // level match = 30
    if (userLevel.HasValue && userLevel.Value == c.TargetLevel) score += 30;
    // field overlap = up to 40, country overlap = up to 25
    score += OverlapScore(preferredFields, fosFields, maxPoints: 40);
    score += OverlapScore(NormalizeCountries(preferredCountries),
                          NormalizeCountries(countries), maxPoints: 25);
    // funding bonus = 5 for large awards
    if (c.FundingAmountUsd >= 10_000) score += 5;

    scored.Add((c.Id, Math.Min(score, 100), c.TitleEn, c.TitleAr));
}
var top = scored.OrderByDescending(x => x.Score).Take(topN).ToList();
\end{lstlisting}
\end{english}

هنشوف إيه في الأسطر دي، واحدة واحدة:
\begin{itemize}
\item \code{if (userLevel == c.TargetLevel) score += 30}: تطابق المستوى الأكاديميّ
يعطي ٣٠ نقطة ثابتة — أقوى إشارة مفردة.
\item \code{OverlapScore(preferredFields, ...)}: التقاطع بين مجالات سارة ومجالات
المنحة يعطي حتى ٤٠ نقطة، موزَّعة بحسب عدد المجالات المشتركة.
\item \code{NormalizeCountries(...)}: نوحّد أسماء الدول («مصر» و\en{Egypt} و\en{EG})
إلى مفتاح واحد قبل المقارنة، حتى لا يفشل التطابق بسبب اختلاف الكتابة.
\item \code{Math.Min(score, 100)}: نسقُف النتيجة عند ١٠٠ لتبقى نسبةً مفهومة.
\item \code{OrderByDescending(...).Take(topN)}: نرتّب تنازليًّا ونأخذ الأعلى فقط.
\end{itemize}

\begin{idea}[الفكرة: ذكاء حسابيّ يمكن الدفاع عنه]
لأنّ الحساب صريح (٣٠ + ٤٠ + ٢٥ + ٥)، تستطيع سارة أن تفهم لماذا ظهرت لها منحة بعينها،
وتستطيع اللجنة أن تتحقّق من النتيجة بالورقة والقلم. لا صندوق أسود، ولا هلوسة، ولا
تكلفة استدعاء نموذج خارجيّ. النتيجة تُخزَّن في صفّ \code{AiInteraction} لتتبّع التكلفة
والاستخدام.
\end{idea}

بجانب الترشيح، تضغط سارة على «افحص الأهليّة» في صفحة منحة، فيعمل
\code{CheckEligibilityCommand} ويقارن \code{LocalAiService} ملفها بثلاثة معايير:
المستوى، والدولة، والمجال. ثمّ تجمع الدالة \code{DeriveVerdict} هذه المعايير في حكم
واحد بقاعدة بسيطة: أيّ «لا» صريحة تعني \en{NotEligible}، وأيّ «جزئيّ» أو «غير معروف»
يعني \en{PartiallyEligible}، وكلّها «نعم» يعني \en{Eligible}.

\begin{center}
\begin{LTR}
\begin{tikzpicture}[node distance=0.7cm and 1.0cm]
  \node[sbox,text width=3.0cm] (in) {\small profile + listing criteria};
  \node[dbox,right=1.0cm of in,text width=3.0cm] (d) {\small any criterion = ``no''?};
  \node[fbox,right=1.1cm of d] (no) {\small NotEligible};
  \node[dbox,below=0.7cm of d,text width=3.0cm] (d2) {\small any ``partial'' / ``unknown''?};
  \node[tbox,right=1.1cm of d2] (part) {\small PartiallyEligible};
  \node[srcbox,below=0.7cm of d2,text width=2.6cm] (yes) {\small Eligible};

  \draw[arr] (in) -- (d);
  \draw[arr] (d) -- node[lbl,above] {yes} (no);
  \draw[arr] (d) -- node[lbl,right] {no} (d2);
  \draw[arr] (d2) -- node[lbl,above] {yes} (part);
  \draw[arr] (d2) -- node[lbl,right] {no} (yes);
\end{tikzpicture}
\end{LTR}
\end{center}

\begin{keyp}[نقطة مهمة: الحكم مُحافِظ عمدًا]
لاحظي أنّ «غير المعروف» يُعامَل مثل «الجزئيّ» لا مثل «النعم». لو كان ملف سارة ناقصًا،
لن يعِدَها النظام بأهليّة كاملة زورًا، بل يقول «جزئيّ» ويطلب منها إكمال الملف. هذا
اختيار تصميميّ محافظ: نُفضّل تحفُّظًا صادقًا على تفاؤل كاذب.
\end{keyp}

\subsection{المحطّة الرابعة: اكتشاف منحة وتقديم طلب}

تختار سارة منحة وتبدأ طلبًا. الطلب يعيش أوّلًا كمسوَّدة (\en{Draft})، تملأ فيها
النموذج وترفع المستندات. حين تضغط «قدِّم»، ترسل الواجهة \code{SubmitApplicationCommand}
فينفّذه \code{SubmitApplicationCommandHandler}. لا يكتفي المعالج بتغيير الحالة، بل
يحرس عدّة قواعد: يتحقّق من ملكيّة الطلب، ويمنع التقديم بعد انتهاء الموعد النهائيّ،
ويتأكّد أنّ كل مستند مرفق يعود فعلًا لملف رفعته سارة عبر خزنة المستندات المفحوصة.

\begin{english}
\begin{lstlisting}[language={[Sharp]C}]
// state machine: Draft -> Pending only if the transition is legal
ApplicationStateMachine.EnsureTransition(oldStatus, ApplicationStatus.Pending);
// deadline guard, even for a Draft that sat past the deadline
if (application.Scholarship is { } sch && sch.Deadline < DateTimeOffset.UtcNow)
    throw new ConflictException("This scholarship's deadline has passed.");

application.Status = ApplicationStatus.Pending;
application.SubmittedAt = DateTimeOffset.UtcNow;
application.RaiseDomainEvent(new ApplicationSubmittedEvent(
    application.Id, application.StudentId, application.ScholarshipId));
await db.SaveChangesAsync(ct);
\end{lstlisting}
\end{english}

الفكرة في الأسطر دي:
\begin{itemize}
\item \code{EnsureTransition}: آلة حالة صغيرة ترفض أيّ انتقال غير قانونيّ (مثلًا من
مقبول إلى مسوَّدة) — القاعدة في مكان واحد لا مبعثرة.
\item حارس الموعد النهائيّ: يمنع تقديم مسوَّدة قديمة بعد فوات الأوان.
\item \code{RaiseDomainEvent(new ApplicationSubmittedEvent(...))}: يسجّل الحدث على
الكيان، ثمّ يُنشَر بعد \code{SaveChangesAsync} فيوقظ الجهة المانحة.
\end{itemize}

\subsection{المحطّة الخامسة: الجهة المانحة تراجع وتردّ (أوّل ردّ فعل)}

الآن يدخل دورٌ آخر للمرّة الأولى. حين قدّمت سارة طلبها، وصل إشعار للجهة المانحة
صاحبة المنحة. تفتح الجهة الطلب وتقرّر: قبول، أو وضعه تحت المراجعة، أو رفض بسبب.
القرار يمرّ عبر \code{ReviewApplicationCommand} الذي ينفّذه المعالج التالي —
لاحظي حارس الملكيّة والإشعار الذي يُرسَل لسارة في نهايته:

\begin{english}
\begin{lstlisting}[language={[Sharp]C}]
// only the scholarship's owning provider may review its applications
if (application.Scholarship.OwnerScholarshipProviderId != currentUser.UserId)
    throw new ForbiddenAccessException();

ApplicationStateMachine.EnsureTransition(oldStatus, request.Status);
application.Status = request.Status;
application.DecisionAt = DateTimeOffset.UtcNow;
application.DecisionReason = request.DecisionReason;   // e.g. reject reason

await db.SaveChangesAsync(ct);
await notifications.DispatchAsync(
    application.StudentId, NotificationType.ApplicationStatusChanged,
    new NotificationParams { StatusText = request.Status.ToString(),
                             Reason = request.DecisionReason }, null, null, ct);
\end{lstlisting}
\end{english}

هنشوف الردّ ماشي إزّاي:
\begin{itemize}
\item \code{OwnerScholarshipProviderId != currentUser.UserId}: حارسٌ يمنع أيّ جهة من
لمس طلبات منحة لا تملكها — عزلٌ صارم بين الجهات.
\item \code{DecisionReason}: عند الرفض تُلزَم الجهة بسببٍ نصّيّ، فتراه سارة صراحةً بدل
رفضٍ صامت — قاعدة عمل حقيقيّة في النظام.
\item \code{notifications.DispatchAsync(application.StudentId, ...)}: يُوجَّه الإشعار
إلى سارة تحديدًا، فيظهر في جرسها لحظيًّا عبر \en{SignalR} مع نصّ الحالة والسبب.
\end{itemize}

\begin{spotlight}[ركّز هنا لإبهار اللجنة: كل خطوة لها ردّ فعل من دورٍ آخر]
هذه أوّل نقطة تُظهر فيها الرحلةُ جوهرَها: فِعلُ سارة (تقديم الطلب) أطلق حدثًا،
والحدثُ أيقظ الجهة المانحة، وردُّ الجهة أطلق إشعارًا رجع إلى سارة. أربعة أدوار
(طالب، وجهة، وأدمن، واستشاريّ) لا تتحدّث مباشرةً، بل تتحاور عبر أحداث وإشعارات
مصمَّمة. أظهر للجنة هذا «التعاون غير المتزامن»: تكاملٌ بلا اقتران متين.
\end{spotlight}

\subsection{المحطّة السادسة: حجز استشاريّ ودفع Stripe (الحجز يُجمّد المبلغ)}

تريد سارة مشورةً بشريّة، فتحجز جلسة مع استشاريّ. عند الحجز يعمل
\code{RequestBookingCommand}. هنا تظهر لمسة الدفع الذكيّة: النظام لا يسحب المال فورًا،
بل يُنشئ نيّة دفع \en{Stripe} بأسلوب القبض اليدويّ
(\code{captureMethod: "manual"}). هذا يعني أنّ بطاقة سارة تُصرَّح (\en{authorized})
والمبلغ يُجمَّد، لكنّه لا يُؤخَذ بعد. يبقى الحجز في حالة \en{Requested} وينتظر ردّ
الاستشاريّ، ويصل إشعار للاستشاريّ بأنّ طلبًا ينتظره.

\begin{why}
لماذا نجمّد المبلغ بدل سحبه؟ لأنّ الاستشاريّ قد يرفض أو لا يردّ. لو سحبنا المال
فورًا لاضطررنا لاستردادٍ لاحقٍ معقّد. التجميد (\en{authorize}) يحجز المبلغ دون أخذه،
فإن قبل الاستشاريّ نقبض، وإن رفض نُفرِج عن التجميد بلا حركة ماليّة فعليّة. المال
«في الضمان» (\en{escrow}) حتى تُحسَم الجلسة.
\end{why}

\subsection{المحطّة السابعة: الاستشاريّ يقبل — نتبّع الأمر من أوّله لآخره}

هذه أهمّ محطّة تقنيًّا، فلنتتبّع أمرًا واحدًا كاملًا من الواجهة إلى الإشعار. حين يضغط
الاستشاريّ «اقبل»، يعمل \code{AcceptBookingCommand}. أوّلًا تعريف الأمر نفسه — سطرٌ
واحد يحمل مُعرِّف الحجز، ومزيّنٌ بسِمة التدقيق:

\begin{english}
\begin{lstlisting}[language={[Sharp]C}]
[Auditable(AuditAction.BookingAccepted, "ConsultantBooking",
    TargetIdProperty = nameof(BookingId),
    SummaryTemplate = "Booking accepted: {TargetId}")]
public sealed record AcceptBookingCommand(Guid BookingId) : IRequest;
\end{lstlisting}
\end{english}

سِمة \code{[Auditable]} تجعل \en{MediatR} يكتب صفَّ تدقيق تلقائيًّا قبل تنفيذ المعالج،
بلا سطر واحد داخل المنطق. الآن قلب المعالج \code{AcceptBookingCommandHandler}:

\begin{english}
\begin{lstlisting}[language={[Sharp]C}]
// 1. authorization: only the owning consultant, only a Requested booking
if (!_currentUser.IsInRole("Consultant")) throw new UnauthorizedAccessException(...);
if (booking.ConsultantId != consultantId) throw new UnauthorizedAccessException(...);
if (booking.Status != BookingStatus.Requested) throw new BookingDomainException(...);

// 2. capture the money that was only held at request time
var captureResult = await _stripeService.CapturePaymentIntentAsync(
    booking.StripePaymentIntentId!, amountCents, idempotencyKey, ct);
payment.Status = PaymentStatus.Captured;   // + platform / payee split

// 3. flip state and provision the ACS video room
booking.Status = BookingStatus.Confirmed;
var room = await _meetingService.CreateRoomAsync(booking.Id, ct);
booking.MeetingRoomId = room.RoomId;

// 4. persist, THEN publish the domain event
await _context.SaveChangesAsync(cancellationToken);
await _publisher.Publish(
    new BookingConfirmedEvent(booking.Id, booking.StudentId, booking.ConsultantId), ct);
\end{lstlisting}
\end{english}

هنشوف المعالج ماشي إزّاي، مرحلة مرحلة:
\begin{itemize}
\item التصريح: ثلاث حراسات تتأكّد أنّ المُنادِي استشاريّ، وأنّه صاحب الحجز، وأنّ
الحجز ما زال \en{Requested} — لا يُقبَل حجزٌ مرّتين.
\item \code{CapturePaymentIntentAsync}: هنا يتحوّل التجميد إلى قبضٍ فعليّ. المبلغ
الذي جُمِّد في المحطّة السابقة يُؤخَذ الآن، وتُحدَّث حصّة المنصّة وحصّة الاستشاريّ.
\item \code{CreateRoomAsync}: تُهيَّأ غرفة \en{ACS} للجلسة ويُحفَظ \code{MeetingRoomId}
على الحجز — نفس الغرفة الكسولة التي شرحناها سابقًا.
\item \code{SaveChangesAsync} ثمّ \code{Publish}: أوّلًا نثبّت الحقائق في قاعدة البيانات،
وبعد نجاح الحفظ فقط نُطلق \code{BookingConfirmedEvent}. الترتيب مقصود: لا نُشعِر أحدًا
بشيء لم يُحفَظ.
\end{itemize}

هذا هو تدفّق الأمر الواحد مرسومًا — من الواجهة حتى الجرس:

\begin{center}
\begin{LTR}
\begin{tikzpicture}[node distance=0.5cm and 0.55cm]
  \node[actor,text width=2.2cm] (fe) {React\\frontend};
  \node[pbox,right=0.55cm of fe,text width=2.7cm] (ctrl) {BookingsController\\.Accept};
  \node[gbox,right=0.55cm of ctrl,text width=2.4cm] (med) {IMediator\\.Send};
  \node[pbox,right=0.55cm of med,text width=2.7cm] (hdlr) {AcceptBooking\\Handler};

  \node[tbox,below=0.75cm of hdlr,text width=2.7cm] (stripe) {Stripe capture\\+ ACS room};
  \node[dbbox,left=0.6cm of stripe,text width=2.2cm] (db) {SQL Server\\SaveChanges};
  \node[tbox,left=0.55cm of db,text width=2.6cm] (evt) {Publish\\BookingConfirmedEvent};
  \node[sbox,left=0.55cm of evt,text width=2.5cm] (notif) {NotificationHub\\bell -> student};

  \draw[arr] (fe) -- node[lbl,above]{POST} (ctrl);
  \draw[arr] (ctrl) -- node[lbl,above]{1 line} (med);
  \draw[arr] (med) -- (hdlr);
  \draw[arr] (hdlr) -- (stripe);
  \draw[arr] (stripe) -- (db);
  \draw[arr] (db) -- node[lbl,above]{then} (evt);
  \draw[arr] (evt) -- (notif);
  \draw[arrd] (notif.south) to[out=270,in=270] node[lbl,below]{push} (fe.south);
\end{tikzpicture}
\end{LTR}
\end{center}

من يلتقط \code{BookingConfirmedEvent}؟ معالجُ حدثٍ منفصل تمامًا،
\code{BookingConfirmedEventHandler}، لا يعرف شيئًا عن \en{Stripe} ولا عن الغرفة —
مهمّته الوحيدة إشعار سارة أنّ حجزها تأكّد:

\begin{english}
\begin{lstlisting}[language={[Sharp]C}]
public async Task Handle(BookingConfirmedEvent notification, CancellationToken ct)
{
    var view = await BookingNotificationContext.LoadAsync(db, notification.BookingId, ct);
    if (view is null) { /* log + return */ return; }

    await notifications.DispatchAsync(
        notification.StudentId,                       // recipient = the student
        NotificationType.BookingConfirmed,
        new NotificationParams { CounterpartyName = view.ConsultantName,
            StartAtText = BookingNotificationContext.FormatStartAt(view.ScheduledStartAt) },
        deepLink: $"/student/bookings/{notification.BookingId}",
        idempotencyKey: $"booking-confirmed:{notification.BookingId:N}", ct);
}
\end{lstlisting}
\end{english}

النقاط المهمّة في ردّ الفعل ده:
\begin{itemize}
\item \code{DispatchAsync(notification.StudentId, ...)}: الإشعار يُوجَّه لسارة تحديدًا،
فيصل جرسها لحظيًّا عبر \en{SignalR}.
\item \code{deepLink}: رابطٌ عميق يفتح صفحة الحجز مباشرة عند الضغط — لا بحث يدويّ.
\item \code{idempotencyKey}: مفتاحٌ فريد يمنع تكرار نفس الإشعار لو أُعيد إطلاق الحدث.
\item فصلٌ نظيف: معالج القبول لا يعرف كيف يُشعِر أحدًا، ومعالج الإشعار لا يعرف شيئًا
عن الدفع. كلٌّ منهما مسؤوليّة واحدة، والحدث هو الجسر بينهما.
\end{itemize}

\begin{spotlight}[ركّز هنا لإبهار اللجنة: رحلة واحدة تلمس كل أجزاء النظام]
تتبّعنا أمرًا واحدًا لمس كلَّ طبقة وخدمة تقريبًا: الواجهة، والمتحكّم، و\en{MediatR}،
والمعالج، و\en{Stripe} (دفع)، و\en{ACS} (فيديو)، وقاعدة البيانات، والحدث الدوميّ،
و\en{SignalR} (إشعار)، وسجلّ التدقيق (\en{[Auditable]}). هذا هو التماسك المعماريّ:
كل قطعة في مكانها، ومسارٌ واحد يشبك بينها جميعًا بلا فوضى. أظهر للجنة هذه اللقطة —
إنّها ملخّص النظام كلّه في أمرٍ واحد.
\end{spotlight}

\subsection{المحطّتان الثامنة والتاسعة: الجلسة والتقييم}

في موعد الجلسة يضغط الطرفان «انضمّ» ضمن نافذة \en{$\pm$15} دقيقة. يعمل
\code{RecordMeetingJoinCommand} فيُصدِر رمز وصول \en{ACS} قصير العمر لكل مشارك،
وتجري المكالمة بين المتصفّحَين مباشرةً بينما يبقى خادمنا مجرّد منسّق. عند الاتصال
يبدأ التسجيل تلقائيًّا ويُخزَّن في \en{Blob Storage}. النظام يسجّل أيضًا وقت انضمام
كل طرف، وهذه الإشارة يقرؤها لاحقًا مسحُ الغياب الآليّ (\en{no-show}) عبر \en{Hangfire}.

بعد الجلسة تقيّم سارة الاستشاريّ عبر \code{SubmitConsultantRatingCommand}: نجوم من
واحد إلى خمسة مع تعليق اختياريّ. التقييم يُغذّي متوسّط الاستشاريّ الظاهر لكل الطلاب
مستقبلًا، وبذلك يردّ فعلُ سارة الصغير على تجربة طلابٍ لم تلتقِهم بعد.

\subsection{المحطّة العاشرة: منشور في المجتمع وردّ حيّ}

تطرح سارة سؤالًا في المجتمع، فيُنشَأ منشور ثنائيّ اللغة. طالبٌ آخر متّصلٌ الآن يرى
الردّ يظهر لحظيًّا دون تحديث الصفحة، لأنّ \en{CommunityHub} يدفع التحديث عبر
\en{WebSocket}. هنا يظهر دورٌ آخر من الأدوار: لو خالف منشورٌ القواعد، يستطيع الأدمن
إخفاءه أو الردّ على بلاغ — نفس نمط الحدث والإشعار الذي رأيناه، لكن في سياق الإشراف.

\begin{note}
لاحظي أنّ الفكرة نفسها تتكرّر في كل محطّة: أمرٌ يغيّر الحالة، ثمّ حدثٌ يُنشَر، ثمّ
دفعٌ لحظيّ عبر \en{SignalR} لمن يعنيه الأمر. تعلّمُ نمطٍ واحد يفتح لك النظام كلّه.
\end{note}

\subsection{المحطّة الحادية عشرة: سؤال لمساعد الذكاء الاصطناعيّ بمصادر}

تختم سارة يومها بسؤال مساعد الشات: «ما الأوراق المطلوبة عادةً؟». يعمل أمر الشات،
وينادي \code{LocalAiService.AskAsync}. هنا يتجلّى \en{RAG} (التوليد المُعزَّز
بالاسترجاع): أوّلًا يسترجع المنفذ \code{IKnowledgeRetriever} أقرب الوثائق من قاعدة
المعرفة، ثمّ يُبقي فقط ما تجاوزت درجتُه العتبة \code{RagMinScore}، ثمّ يبني الإجابة
من أعلى وثيقة مطابقة — ويُرفِق مصادرها.

\begin{english}
\begin{lstlisting}[language={[Sharp]C}]
var docs = await retriever.RetrieveAsync(msg, opts.Value.RagTopK, ct);
var relevant = docs.Where(d => d.Score >= opts.Value.RagMinScore).ToList();

string reply = relevant.Count > 0
    ? $"{lead}\n\n{(arabic ? relevant[0].ContentAr : relevant[0].ContentEn).Trim()}"
    : RouteChat(msg);            // keyword fallback when nothing scores high enough

return new AiChatResponse(Message: reply, /* tokens, cost */
    Sources: RagSupport.ToSources(relevant, arabic));   // cited sources
\end{lstlisting}
\end{english}

المهمّ في الأسطر دي:
\begin{itemize}
\item \code{RetrieveAsync}: يبحث في قاعدة المعرفة عن أقرب \code{RagTopK} وثيقة للسؤال.
\item \code{d.Score >= RagMinScore}: عتبةُ ثقة — لو لم تتجاوز أيّ وثيقة العتبة، لا نخترع
إجابة، بل نرجع لموجِّه كلماتٍ مفتاحيّة (\code{RouteChat}) يعطي نصيحة عامّة آمنة.
\item \code{Sources: RagSupport.ToSources(...)}: تُرفَق المصادر مع الردّ، فترى سارة
من أين جاءت الإجابة — لا صندوق أسود.
\end{itemize}

\begin{keyp}[نقطة مهمة: نفس التصميم يعمل بلا مفتاح خارجيّ]
هذا المسار «المحلّيّ» يشغّل \en{RAG} كاملًا بلا استدعاء \en{GPT} إطلاقًا. لو ضُبِط
مزوّد \en{Azure OpenAI}، يُحقَن محوّلٌ آخر يولّد إجابةً حرّة فوق نفس الوثائق المسترجَعة.
تبديل المزوّد قرارُ إعداد لا قرار كود — نفس نمط المنافذ في كل النظام.
\end{keyp}

\subsection{من فعل ماذا: خريطة الأدوار في الرحلة}

لنجمع الرحلة في جدول واحد يبيّن الأمر عند كل خطوة والدور الذي يردّ:

\begin{center}
\begin{tabular}{L{3.4cm} L{4.3cm} L{4.3cm}}
\toprule
المرحلة & الأمر / المعالج & الدور الذي يردّ \\
\midrule
التسجيل & \en{Identity} + تأكيد البريد & لا أحد (بوّابة دخول) \\
التوصية والأهليّة & \en{GenerateRecommendations} & \en{LocalAiService} (حسابيّ) \\
تقديم الطلب & \en{SubmitApplicationCommand} & الجهة المانحة تُشعَر \\
مراجعة الطلب & \en{ReviewApplicationCommand} & سارة تُشعَر بالنتيجة والسبب \\
حجز الجلسة & \en{RequestBookingCommand} & الاستشاريّ يُشعَر (المبلغ مُجمَّد) \\
قبول الحجز & \en{AcceptBookingCommand} & سارة تُشعَر (القبض يتمّ) \\
المجتمع والإشراف & منشور / بلاغ & طالب آخر أو أدمن يردّ \\
\bottomrule
\end{tabular}
\end{center}

\begin{qa}[سؤال متوقّع: هل الأدوار مقترنة اقترانًا متينًا؟]
لا. لاحظ أنّ أيّ معالج لا ينادي دورًا آخر مباشرةً. معالج القبول يُطلق حدثًا فقط،
ولا يعرف أنّ هناك معالج إشعار يستمع. لو حذفنا الإشعار غدًا، لن يتغيّر سطرٌ في منطق
القبول. هذا اقتران رخو (\en{loose coupling}) عبر الأحداث: الأدوار تتعاون دون أن
يعرف بعضُها تفاصيل بعض.
\end{qa}

\begin{summary}[الخلاصة في سطور]
\begin{itemize}
\item تابعنا سارة عبر إحدى عشرة محطّة، فلمسنا كل جزء من النظام في قصّة واحدة متّصلة.
\item كل خطوة أمرٌ (\en{command}) يغيّر الحالة، وكثيرٌ منها يُطلق حدثًا دوميًّا بعد
\code{SaveChangesAsync}، والحدثُ يوقظ دورًا آخر عبر إشعار لحظيّ.
\item تتبّعنا \code{AcceptBookingCommand} كاملًا: واجهة، فمتحكّم، فمعالج، فقبضُ \en{Stripe}،
فغرفة \en{ACS}، فحفظٌ، فحدثٌ، فإشعارٌ للجرس — رحلة أمرٍ واحدة تلخّص المعمار كلّه.
\item الذكاء الاصطناعيّ ظهر ثلاث مرّات: توصيةٌ حسابيّة قابلة للتفسير (٣٠+٤٠+٢٥+٥)،
وفحصُ أهليّة بحكمٍ محافظ، وشاتٌ يردّ بإجابةٍ مع مصادرها عبر \en{RAG}.
\item الأدوار الأربعة (طالب، وجهة، وأدمن، واستشاريّ) تتحاور عبر أحداث وإشعارات لا عبر
نداءات مباشرة — تكاملٌ مصمَّم باقترانٍ رخو، وهو سرّ تماسك النظام.
\end{itemize}
\end{summary}
